\documentclass[10pt]{article}
\providecommand{\FIGDIR}{fig}
\newcommand{\DIGESTDATE}{6 August 2026}
% ===================== PM Research Watch - shared preamble =====================
\usepackage[T1]{fontenc}
\usepackage[utf8]{inputenc}
\usepackage{amsmath,amssymb}
\usepackage{tgpagella}
\usepackage{tgheros}
\renewcommand{\sfdefault}{qhv}
\renewcommand{\ttdefault}{lmtt}
\usepackage{microtype}
\usepackage[english]{babel}

\usepackage[a4paper,top=17mm,bottom=17mm,left=17mm,right=17mm,
            headsep=6mm,footskip=10mm]{geometry}
\usepackage{graphicx}
\usepackage{xcolor}
\usepackage{colortbl}
\usepackage{tikz}
\usetikzlibrary{calc}
\usepackage[most]{tcolorbox}
\usepackage{titlesec}
\usepackage{fancyhdr}
\usepackage{booktabs}
\usepackage{array}
\usepackage{longtable}
\usepackage{enumitem}
\usepackage{multicol}
\usepackage{ragged2e}
\usepackage[hidelinks]{hyperref}

% ------------------------------------------------------------------ palette
\definecolor{Ink}    {HTML}{14262E}
\definecolor{Deep}   {HTML}{0F4C5C}
\definecolor{Teal}   {HTML}{1F7A8C}
\definecolor{Sky}    {HTML}{4FA3C4}
\definecolor{Sage}   {HTML}{6FA287}
\definecolor{Amber}  {HTML}{E3A008}
\definecolor{Clay}   {HTML}{B5643C}
\definecolor{Coral}  {HTML}{D1495B}
\definecolor{Violet} {HTML}{6A5B9E}
\definecolor{Slate}  {HTML}{4A6572}
\definecolor{Paper}  {HTML}{FBF9F5}
\definecolor{Mist}   {HTML}{EEEAE2}
\definecolor{Rule}   {HTML}{D9D3C9}

\pagecolor{Paper}
\color{Ink}

\hypersetup{colorlinks=true, linkcolor=Deep, urlcolor=Teal, citecolor=Teal,
            pdftitle={PM Research Watch}, pdfauthor={Automated daily digest}}

% ------------------------------------------------------------------ headings
\titleformat{\section}
  {\sffamily\bfseries\large\color{Deep}}{}{0pt}
  {}[{\vspace{1.2mm}\color{Amber}\titlerule[1.1pt]}]
\titlespacing*{\section}{0pt}{5.5mm}{2.8mm}

\titleformat{\subsection}
  {\sffamily\bfseries\normalsize\color{Teal}}{}{0pt}{}
\titlespacing*{\subsection}{0pt}{3.6mm}{1.4mm}

\setlist[itemize]{leftmargin=4.2mm, itemsep=1.1mm, topsep=1.3mm, parsep=0pt,
                  label=\textcolor{Teal}{\footnotesize$\blacktriangleright$}}

% ------------------------------------------------------------------ header/footer
\pagestyle{fancy}
\fancyhf{}
\renewcommand{\headrulewidth}{0pt}
\renewcommand{\footrulewidth}{0pt}
\fancyhead[L]{\sffamily\scriptsize\color{Slate}PM RESEARCH WATCH\;\textcolor{Rule}{\textbar}\;\DIGESTDATE}
\fancyhead[R]{\sffamily\scriptsize\color{Slate}Particulate matter monitoring \& health effects}
\fancyfoot[C]{\sffamily\scriptsize\color{Slate}\thepage}

% ------------------------------------------------------------------ components
% metric tile
\newcommand{\metric}[3]{%
  \begin{tcolorbox}[enhanced, colback=white, colframe=Rule, boxrule=0.5pt,
    arc=1.6mm, left=2.2mm, right=2.2mm, top=1.8mm, bottom=1.8mm,
    borderline west={1.7mm}{0pt}{#3}, nobeforeafter, valign=center,
    height=17mm, width=\linewidth]
    \RaggedRight
    {\sffamily\bfseries\fontsize{17}{17}\selectfont\color{#3}#1}\\[0.5mm]
    {\sffamily\fontsize{6.5}{7.4}\selectfont\color{Slate}#2}
  \end{tcolorbox}}

% pull-quote / key insight box
\newtcolorbox{keybox}[1][]{enhanced, colback=Mist, colframe=Mist,
  arc=1.6mm, left=3.2mm, right=3.2mm, top=2.4mm, bottom=2.4mm,
  borderline west={1.7mm}{0pt}{Amber}, #1}

% subtopic band
\newcommand{\band}[2]{%
  \vspace{3.2mm}\noindent
  \begin{tcolorbox}[enhanced, colback=#2, colframe=#2, arc=1.2mm,
    left=2.6mm, right=2.6mm, top=1.1mm, bottom=1.1mm, boxrule=0pt,
    width=\linewidth, nobeforeafter]
    {\sffamily\bfseries\small\color{white}#1}
  \end{tcolorbox}\vspace{1.6mm}}

% one paper entry: {short cite}{title}{journal}{doi}{body}
\newcommand{\paperentry}[5]{%
  \par\vspace{1.9mm}\noindent
  {\sffamily\bfseries\footnotesize\color{Ink}#1}\;%
  {\sffamily\fontsize{7}{8}\selectfont\color{Slate}#3}\\[0.4mm]
  {\footnotesize\itshape\color{Deep}#2}\\[0.6mm]
  {\footnotesize #4}\\[0.3mm]
  {\fontsize{6.6}{7.6}\selectfont\color{Slate}\texttt{doi:}\,\href{https://doi.org/#5}{\texttt{\detokenize{#5}}}}
  \par}

% caption style
\newcommand{\figcap}[1]{\par\vspace{1.0mm}
  {\fontsize{7.2}{8.6}\selectfont\color{Slate}\textsf{\textbf{Fig.}}\; #1\par}\vspace{1.5mm}}

\providecommand{\permil}{\textperthousand}

\setlength{\parindent}{0pt}
\setlength{\parskip}{0pt}
\renewcommand{\arraystretch}{1.18}


\begin{document}

% ==================================================================== MASTHEAD
\thispagestyle{empty}
\begin{tikzpicture}[remember picture, overlay]
  \fill[Deep] (current page.north west) rectangle
      ($(current page.north east)+(0,-67mm)$);
  \fill[Amber] ($(current page.north west)+(0,-67mm)$) rectangle
      ($(current page.north east)+(0,-68.4mm)$);
  \foreach \x/\y/\r in {22/12/0.9, 41/26/1.5, 63/9/0.7, 88/31/1.1, 112/17/1.9,
                        138/29/0.8, 159/13/1.3, 178/54/1.0, 196/21/1.6,
                        30/44/0.6, 74/57/1.2, 125/48/0.9, 168/6/0.8, 191/59/0.7,
                        52/18/1.0, 100/60/0.6, 148/20/1.1, 66/49/1.4, 205/45/0.8}
    \fill[Sky, opacity=0.30] ($(current page.north west)+(\x mm,-\y mm)$) circle (\r mm);
\end{tikzpicture}

\vspace*{4mm}
{\sffamily\fontsize{7.6}{9}\selectfont\color{Sky}%
 AUTOMATED DAILY LITERATURE DIGEST\;\textbullet\;PhD RESEARCH WATCH\par}
\vspace{2.2mm}
{\sffamily\bfseries\fontsize{27}{29}\selectfont\color{white}%
 Particulate Matter\par}
{\sffamily\fontsize{27}{29}\selectfont\color{Sky}%
 Monitoring \& Health Effects\par}
\vspace{5.0mm}
{\sffamily\fontsize{8.4}{10}\selectfont\color{white}%
 Issue 2026-08-06\;\;\textcolor{Amber}{\textbar}\;\;Entry window 6 August 2026\;\;%
 \textcolor{Amber}{\textbar}\;\;27 records in scope\par}
\vspace{18mm}

% -------------------------------------------------------------- metric strip
\noindent
\begin{minipage}[t]{0.190\linewidth}\metric{27}{RECORDS IN SCOPE\\(23 rejected, 9 held)}{Deep}\end{minipage}\hfill
\begin{minipage}[t]{0.190\linewidth}\metric{7}{SENSING RECORDS\\(26\% OF CORPUS)}{Teal}\end{minipage}\hfill
\begin{minipage}[t]{0.190\linewidth}\metric{8}{SUBTOPIC CLUSTERS\\REPRESENTED}{Sky}\end{minipage}\hfill
\begin{minipage}[t]{0.190\linewidth}\metric{9}{TIER-A\\RECORDS}{Amber}\end{minipage}\hfill
\begin{minipage}[t]{0.190\linewidth}\metric{6}{EFFECT ESTIMATES\\EXTRACTED}{Clay}\end{minipage}

\vspace{6mm}

% -------------------------------------------------------------- signal of day
\section*{\sffamily\bfseries\color{Deep}Signal of the day}
\vspace{-1.5mm}{\color{Amber}\rule{\linewidth}{1.1pt}}\vspace{2.5mm}

\begin{keybox}
{\sffamily\bfseries\small\color{Deep}What people \emph{smell} moves protective behaviour
roughly five times more than what the network \emph{measures}. The last metre of a
monitoring system is perceptual, and it is currently unengineered.}\\[1.4mm]
\footnotesize
Slavik et al.\ (\textit{Environ Challenges}) linked survey responses from
\textbf{2{,}086 parents} in wildfire-prone western North America to three and a half
years of smoke-monitoring data (2020--2023). Parents who \emph{perceived} smoke exposure
took \textbf{1.11 more protective actions} (95\% CI 0.92--1.30) than those who did not,
whereas each \textbf{10\,$\mu$g/m$^3$ of measured PM$_{2.5}$} on smoke days bought only
\textbf{0.23 additional actions} (0.06--0.40) --- a $\sim$4.8-fold gap in favour of the
subjective channel. The measured-exposure response is also non-monotonic: actions rise,
plateau at moderate concentrations, then rise again. Since sub-visible smoke is
demonstrably harmful, a network that reports accurate concentrations into a population
that acts on sight and smell is delivering a fraction of its achievable health benefit.
Read against the two records in today's corpus that show \emph{mass} is the wrong
currency for the hazard --- Yadav et al.'s Agra year-long series, where PM$_{10}$ mass and
endotoxin are effectively uncorrelated ($r_s=0.10$, $p=0.75$) across an 8-fold seasonal
endotoxin swing, and Lei et al.'s eight-year speciated PM$_1$ record in Dublin --- the
day's through-line is that the two ends of the measurement chain are both weakly coupled
to the thing that matters. Composition is not mass, and communicated risk is not
measured risk.
\end{keybox}

\vspace{3mm}

% Exactly two lines of the last bullet spilled onto an otherwise blank p.2 and two rounds
% of prose trimming did not converge (the paragraph reflows and re-spills). 2 baselines is
% the smallest dose that absorbs it; the 2026-08-04 run log records that 4-7 baselines
% start printing text over the folio, so re-proof the bottom margin if this is ever raised.
\enlargethispage{2\baselineskip}
% -------------------------------------------------------------- exec brief
\section*{\sffamily\bfseries\color{Deep}What changed today}
\vspace{-1.5mm}{\color{Amber}\rule{\linewidth}{1.1pt}}\vspace{2.5mm}

\footnotesize
\begin{itemize}
\item \textbf{The confidence-interval drought broke.} Jung et al.\ (\textit{Ann Rheum
Dis}) report the first peer-reviewed PM ratio estimate with an interval this watch has
seen in five issues: \textbf{OR 1.113 (95\% CI 1.017--1.218)} for rheumatoid arthritis
flare per monthly mean PM$_{2.5}$, across 12{,}583 outpatient visits from 1{,}070
patients, with a within-patient case-crossover sensitivity arm that reproduces the
association at $>$2 weeks of cumulative exposure. PM$_{2.5}$ was the only one of six
pollutants to survive. Six effect estimates carry intervals today against one on 5
August.

\item \textbf{Two well-powered nulls, and both are informative rather than empty.}
Sewor et al.\ found \emph{no} association between 48-hour personal PM$_{2.5}$ and blood
pressure in 622 rural Rwandan children --- at a \textbf{median exposure of
192.9\,$\mu$g/m$^3$}, roughly 40$\times$ the WHO annual guideline. The authors' own
reading is the right one: the exposure--response may simply be flat that far up the
global continuum, which is a statement about where monitoring effort buys epidemiological
information. Separately, Lees et al.\ (preprint) report a PM$_{2.5}$-component mixture
IRR of \textbf{1.31 (0.99--1.73, $p=0.06$)} for hospitalisation in 2{,}511 CKD patients
with \textbf{93\% of the WQS weight on nitrate alone} --- a null whose weight vector is
more interesting than its point estimate.

\item \textbf{Instrumentation side: three papers about the limits of what an instrument
can be asked to mean.} Liu et al.\ (\textit{ACP}) show that inferring the aerosol
mixing-state index $\chi$ from HTDMA $\kappa$ distributions works until externally mixed
\emph{intermediate}-hygroscopicity species violate the method's binary assumption, and
publish size-resolved error bounds so the retrieval can be used honestly. Tian et al.\
(\textit{AMT}) do the analogue for lidar-retrieved CCN profiles with a heterogeneity index
that flags when vertical layering has invalidated the homogeneous-composition assumption.
Yee et al.\ (\textit{AMT}) release UCB-GLOBES, $\sim$27{,}000 GC/EI-MS spectra --- and
report \textbf{$\sim$70\% of them still match nothing}, in their own library or in NIST.

\item \textbf{Imputation gets an honest negative.} Somnugpong's similarity-weighted
seasonal-anchor Stineman method improves on plain Stineman for the UCI dataset under MCAR
and MAR but \emph{does not} beat it on Beijing PM$_{2.5}$. The paper is framed as a
boundary condition rather than a method launch, which is the correct framing and a rare
one --- relevant to anyone reconstructing gappy low-cost sensor series.

\item \textbf{Preclinical cardiac remodelling is now measurable at IQR-scale exposure
contrasts.} Leiser et al.\ (\textit{JACC Adv}) find per-IQR PM$_{2.5}$ associated with
+0.52\% global longitudinal strain (0.20--0.83), $-$1.05\% LVEF per year
($-$1.73 to $-$0.37) and $-$1.39 left atrial volume index ($-$2.25 to $-$0.54) over a
median 4.3 years of serial echocardiography in 1{,}630 people. Note the sign pattern is
internally awkward --- GLS and LAVI move in directions that do not obviously cohere --- and
the paper does not resolve it.

\item \textbf{Sub-Saharan Africa contributed three records, the most in any window so
far} --- Mushore et al.\ (road-edge gradient \emph{reverses sign} between sites),
Makanjuola (563 street sweepers) and Sewor above, all three using handheld or wearable
optical instruments as the primary exposure measurement.
\end{itemize}

\clearpage

% ==================================================================== ANALYTICS
\section*{\sffamily\bfseries\color{Deep}Corpus analytics}
\vspace{-1.5mm}{\color{Amber}\rule{\linewidth}{1.1pt}}\vspace{3mm}

\noindent\includegraphics[width=\linewidth]{\FIGDIR/f1_subtopics.png}
\figcap{Subtopic assignment of the 27 in-scope records, single-label by dominant
contribution. \textit{Sensing, forecasting \& instrumentation} leads with 7 (26\%),
followed by burden/policy (5) and mechanistic toxicology (4). Eight of the ten standing
clusters are populated --- the widest coverage since 27 July --- but \textbf{\textit{Neuro /
mental health} and \textit{Respiratory \& allergic} are both empty as dominant labels},
which is an artefact of assignment rather than of the corpus: four records carry a
respiratory endpoint but sit under sensing or occupational because that is where their
methodological contribution lies.}

\vspace{2mm}
\noindent\includegraphics[width=\linewidth]{\FIGDIR/f2_design_endpoint.png}
\figcap{Left --- study architecture, evenly spread for the first time in a fortnight:
measurement campaigns, cross-sectional observational work and experimental toxicology tie
at 5 each, with modelling/inventory and cohort designs at 4. Right --- health endpoint by
organ system. \textbf{Eleven of 27 records carry no human health endpoint}, consistent
with an instrumentation-heavy window; respiratory (4) leads the endpoints that exist,
ahead of cardiovascular (3) and other-clinical (3). There is no trial or intervention
record in this window at all.}

\vspace{2mm}
\noindent\includegraphics[width=\linewidth]{\FIGDIR/f3_metric_geography.png}
\figcap{Left --- particle metric. Ten records use PM$_{2.5}$ alone and three pair it with
PM$_{10}$. The 11-record unspeciated/emissions bar is genuine heterogeneity rather than
missing detail --- it holds PM$_1$, CCN number, mass-spectral organic aerosol, a
mixing-state index and two records whose exposure is a gas. Two records are resolved to
ultrafine/nanoscale and one to coarse bioaerosol, the latter being the Agra endotoxin
series. Right --- geographic provenance. North America leads (6), then China and the global/multi-region
group (4 each). \textbf{Sub-Saharan Africa contributes 3 and South Asia 3} --- together
22\% of the corpus, which is the strongest representation from the highest-burden regions
this watch has recorded.}

\vspace{2mm}
\noindent\includegraphics[width=\linewidth]{\FIGDIR/f6_lifecourse.png}
\figcap{Life-course windows addressed; counts are non-exclusive. Nine of 27 records carry
a human life-course window. Adults dominate (6). \textbf{The adolescent and older-adult
windows are both empty}, the older-adult gap now running three consecutive issues despite
that group carrying the largest cardiovascular attributable burden. The single in-utero
record is a nitrogen-dioxide study, not a PM study.}

\vspace{2mm}
\noindent\includegraphics[width=\linewidth]{\FIGDIR/f4_heatmap.png}
\figcap{Subtopic $\times$ study-architecture cross-tabulation. The structural feature to
note is that the \textit{Sensing} row is split between measurement campaigns and
modelling with \textbf{no cell in any observational-epidemiology column} --- instrumentation
and health inference remain disjoint within this window, as they have been in every issue
since 27 July. Mechanistic toxicology is again wholly experimental (4/4), with no human
biomarker bridge.}

\vspace{2mm}
\noindent\includegraphics[width=\linewidth]{\FIGDIR/f5_forest.png}
\figcap{All six effect estimates carrying an interval, on a common logarithmic axis.
\textbf{Exposure contrasts are heterogeneous and the estimates are not mutually
comparable}: Jung is per monthly mean PM$_{2.5}$, Li (\textit{Front Pediatr}) is a
categorical $>$75 vs $<$35\,$\mu$g/m$^3$ contrast, Yoon is a short-lag RR per unit
bushfire-smoke PM$_{2.5}$, Lees is per quartile of a weighted-quantile-sum component
index, and the two lowest entries are \textbf{co-pollutant estimates (NO$_2$, O$_3$) and
not PM at all} --- they are plotted so a rollup does not read their absence as pipeline
loss. Read direction and precision, not relative magnitude. Four of the six are
PM estimates, ending a four-issue run in which no peer-reviewed PM ratio estimate with a
confidence interval entered the corpus.}

\clearpage

% ==================================================================== DIGEST
\section*{\sffamily\bfseries\color{Deep}Paper-level digest}
\vspace{-1.5mm}{\color{Amber}\rule{\linewidth}{1.1pt}}\vspace{1mm}

{\fontsize{7.4}{9}\selectfont\color{Slate}Ordered by cluster. Each entry gives the
methodological core and the finding that matters, not an abstract paraphrase.\par}

\band{Sensing, forecasting \& instrumentation\hfill 7 records}{Amber}

\paperentry{Somnugpong 2026}
{Similarity-weighted seasonal-anchor Stineman interpolation for gappy air-quality series}
{Environ Monit Assess}
{SW-SAST keeps a local Stineman anchor and selectively borrows cross-day information when
nearby daily analogues are compatible with the local estimate. Evaluated on UCI Air
Quality and Beijing PM$_{2.5}$ under MCAR, MAR and MNAR at 10--50\% missingness, 30
repeated runs per condition, against a univariate LSTM baseline and with contiguous
gap-length experiments. \textbf{It improves on plain Stineman for UCI, chiefly under MCAR
and MAR, and fails to beat it on Beijing.} The contribution is explicitly a boundary
condition on when seasonal borrowing helps, not a replacement interpolator --- which is
the right claim and an unusually disciplined one. Directly applicable to reconstructing
low-cost-sensor series where dropouts are neither random nor uniform.}
{10.1007/s10661-026-15770-8}

\paperentry{Lei et al.\ 2026}
{Eight-year speciated PM$_1$ record in Dublin: heating emissions versus transboundary transport}
{Atmos Chem Phys}
{Continuous chemically speciated PM$_1$, 2016--2023, resolving two distinct pollution
regimes: short intense events (hours, PM$_1>$100\,$\mu$g/m$^3$) driven by local solid-fuel
home heating, and moderate multi-day events ($<$60\,$\mu$g/m$^3$) from continental
transport. Annual PM$_1$ fell from 6.5 to $<$5.0\,$\mu$g/m$^3$, led by nitrate and
ammonium ($-$0.11 and $-$0.09\,$\mu$g\,m$^{-3}$\,yr$^{-1}$). \textbf{Transported
oxygenated organic aerosol rose (+0.34\,$\mu$g\,m$^{-3}$\,yr$^{-1}$)} while locally formed
OOA barely fell, so the composition shifted even as mass declined --- and 2023 showed a
heating-driven rebound. A demonstration that a mass-trend headline can be exactly wrong
about what a network is now measuring.}
{10.5194/acp-26-10927-2026}

\paperentry{Liu et al.\ 2026}
{From $\kappa$ to $\chi$: particle-resolved evaluation of hygroscopicity-based mixing-state retrieval}
{Atmos Chem Phys}
{PartMC-MOSAIC urban-plume ensembles spanning emissions, ageing and meteorology were used
to compare the particle-resolved reference mixing-state index against the Yuan--Zhao
retrieval from HTDMA $\kappa$ distributions, per diameter over 50--250\,nm. The retrieval
performs well for many populations but \textbf{systematically overestimates $\chi$ when
externally mixed intermediate-hygroscopicity components break its binary
more-/less-hygroscopic assumption}. The useful output is the size-resolved error
distribution, which the authors then apply as uncertainty bounds to long-term HTDMA
records at urban, continental and coastal sites. Model-derived error bars on an
observational retrieval --- the validation is only as good as PartMC's own composition
realism, which the paper does not independently test.}
{10.5194/acp-26-10965-2026}

\paperentry{Yee et al.\ 2026}
{UCB-GLOBES: open mass-spectral library for identified \emph{and} unidentified organic aerosol}
{Atmos Meas Tech}
{$\sim$27{,}000 GC/EI-MS spectra of trimethylsilyl derivatives of semi-volatile organic
aerosol from US and Central Amazon ambient samples plus laboratory SOA, catalogued
\emph{without regard to known chemical identity} and annotated with sample provenance and
predicted properties. Reanalysis of prior ambient data gave an 11\% net increase in
MS-based source categorisation for the southeastern US and lifted Central Amazon matches
from 136 to 375 of 1{,}513 analytes. \textbf{$\sim$70\% of UCB-GLOBES spectra match
nothing} in the library or in NIST --- the honest headline is the size of the unidentified
fraction, not the size of the library.}
{10.5194/amt-19-5027-2026}

\paperentry{Zhao et al.\ 2026}
{WHA air-LSTM: composite morbidity-driven pollution index for next-day outpatient forecasting}
{Biomed Environ Sci}
{A five-level composite index is built \emph{per city} from city-specific
exposure--response relationships across $>$223.7 million hospital visits, then fed with
temperature, humidity and visit history into an LSTM predicting next-day respiratory
outpatient load. Beijing $R^2=0.963$, RMSE 53.5; removing the index degrades RMSE by
44.1\%; performance holds on held-out cities. The index is fitted on morbidity and then
used to predict morbidity, so the incremental-value claim is partly circular --- the
external validation is what carries it. Attractive template for turning a dense sensor
network into an operational early-warning product.}
{10.3967/bes2026.062}

\paperentry{Li et al.\ 2026}
{Interpretable ML for 90-day paediatric allergic-rhinitis exacerbation from clinical and ambient data}
{Front Pediatr}
{Development cohort $n=2{,}000$ linking NHANES (2005--06, 2011--18) to EPA AQS and NOAA
data; external cohort $n=50$ from Wuhan. L2-regularised logistic regression was
prespecified and beat random forest, XGBoost, LightGBM and SVM under nested 5-fold CV.
Internal AUC 0.81 (0.74--0.88), external 0.71 (0.50--0.92). Baseline symptom score,
same-day PM$_{2.5}$ and total IgE dominate SHAP; \textbf{PM$_{2.5}>$75 vs
$<$35\,$\mu$g/m$^3$ gives OR 5.08 (3.45--7.47)}. The external cohort is 50 children on a
different continent --- the authors say plainly this is preliminary transportability, not
external validity, and they are right.}
{10.3389/fped.2026.1865399}

\paperentry{Tian et al.\ 2026}
{A heterogeneity index for when lidar-retrieved CCN vertical profiles stop being trustworthy}
{Atmos Meas Tech}
{ARM ground-based Raman lidar plus surface CCN measurements retrieve vertically resolved
CCN number concentration on the assumption of vertically homogeneous composition. A
$k$-nearest-neighbours classifier on lidar ratio and linear depolarisation ratio infers
dominant aerosol type by altitude, and a derived heterogeneity index quantifies vertical
complexity. Validated against HI-SCALE aircraft in situ data: \textbf{retrieval closure
degrades monotonically as the index rises}, so the metric functions as a self-diagnostic
flag on the retrieval. Single campaign, so the index thresholds are not yet portable.}
{10.5194/amt-19-5051-2026}

\clearpage

\band{Exposure assessment \& modelling\hfill 3 records}{Sky}

\paperentry{Mushore et al.\ 2026}
{Daily outdoor PM$_{2.5}$ and the road-edge effect across The Gambia, Kenya and Mozambique}
{Atmos Environ X}
{Multi-temporal personal exposure measurements (2022--23) interpolated to continuous
surfaces by inverse-distance weighting, validated leave-one-out at $R^2>$0.55, MAE
$<$4\,$\mu$g/m$^3$, RMSE $<$10\,$\mu$g/m$^3$; multi-distance buffers at 10\,m intervals to
1000\,m from major roads. Most sites exceed WHO guidance, The Gambia highest across all
land-cover classes. \textbf{The road-edge gradient reverses sign between sites} ---
PM$_{2.5}$ falls with distance in Rabai and Mozambique and \emph{rises} elsewhere,
implicating wetland, fire and industrial sources that a traffic-proximity exposure proxy
would misattribute entirely. IDW over sparse personal measurements is a weak spatial
model; the buffer result is the transferable part.}
{10.1016/j.aeaoa.2026.100463}

\paperentry{Helbich et al.\ 2026}
{Socioeconomic patterning of the urban external exposome depends on how the exposome is defined}
{Environ Int}
{4{,}640 adults across Bulgaria's five largest cities, with residential NO$_2$, PM, traffic
noise, green space, blue-space distance, impervious surface, walkability and night-time
light. City-specific latent profile analysis was compared against composite indices, then
both were regressed on individual- and area-level SES. High-risk profiles bundle low
green space, high air pollution, noise and imperviousness \emph{with} higher walkability.
\textbf{Higher area-level SES was positively associated with high-risk profiles in most
cities} --- the inverse of the standard environmental-justice expectation --- and composite
indices only partially reproduced the latent-profile findings. The methodological warning
is the finding: operationalisation choice changes the sign of the inequality conclusion.}
{10.1016/j.envint.2026.110419}

\paperentry{Yadav et al.\ 2026}
{Airborne endotoxin on PM$_{10}$ in Agra: an 8-fold seasonal swing that PM mass does not track}
{Research Square (preprint)}
{Full calendar year (Jan--Dec 2022) of ambient endotoxin by chromogenic LAL assay with
concurrent bacterial and fungal bioaerosol, PM$_{10}$ and meteorology in the semi-arid
Indo-Gangetic Basin. Monthly means run 0.92$\pm$0.35 (June) to 7.35$\pm$1.48\,EU\,m$^{-3}$
(December) on winter stagnation. Endotoxin correlates strongly with bacterial
concentration ($r_s=0.93$) and \textbf{not at all with PM$_{10}$ mass ($r_s=0.10$,
$p=0.75$)} --- mineral dust is not the carrier here. USEPA-model hazard quotients stay
well below unity (max 0.0198, children, December). Preprint, single city, LAL assays are
notoriously matrix-sensitive; the decoupling from mass is the durable point.}
{10.21203/rs.3.rs-10523351/v1}

\band{Cardiovascular \& metabolic\hfill 3 records}{Teal}

\paperentry{Leiser et al.\ 2026}
{Ambient PM$_{2.5}$ and preclinical change in cardiac structure and function}
{JACC Adv}
{1{,}630 participants with echocardiography at two visits (2008--11, 2015--18), median 4.3
years apart; residential PM$_{2.5}$ and NO$_2$ from spatio-temporal models as annual means
preceding each visit, adjusted for baseline and time-varying covariates. Per IQR of
PM$_{2.5}$: \textbf{global longitudinal strain +0.52\% (0.20--0.83), LVEF
$-$1.05\%/year ($-$1.73 to $-$0.37), relative wall thickness +0.12 unit (0.05--0.19),
left atrial volume index $-$1.39 ($-$2.25 to $-$0.54)}. NO$_2$ moves the same direction,
smaller. The LAVI direction sits awkwardly with a concentric-remodelling reading and the
paper does not reconcile it; treat the strain and RWT results as the load-bearing ones.}
{10.1016/j.jacadv.2026.103104}

\paperentry{Sewor et al.\ 2026}
{Personal household-air-pollution exposure and blood pressure in rural Rwandan children}
{Environ Res Health}
{622 children aged 8--15 from 626 biomass-cooking households, 48\,h personal PM$_{2.5}$
and black carbon by Ultrasonic Personal Aerosol Sampler, then standardised BP measurement.
Median 48\,h PM$_{2.5}$ \textbf{192.9\,$\mu$g/m$^3$} (IQR 119.7--336.0), BC
9.85\,$\mu$g/m$^3$. \textbf{No association with SBP or DBP} (PM$_{2.5}$--SBP 0.51\,mmHg
per 211.7\,$\mu$g/m$^3$ IQR, 95\% CI $-$0.46 to 1.49) and no clear effect modification by
age or sex. Yet 24\% of the children were hypertensive or had elevated BP against US
reference percentiles. Cross-sectional, and the exposure range sits entirely above the
region of the global exposure--response curve where slope is expected --- a null that
constrains where measurement effort is epidemiologically productive.}
{10.1088/2752-5309/ae8ff9}

\paperentry{Liu et al.\ 2026}
{Long-term ozone and cardiovascular disease: exposure windows, sex and region}
{China CDC Wkly}
{Meta-analysis of cohort studies from PubMed, Web of Science and Embase, stratified in the
way most prior reviews were not. Per 10\,$\mu$g/m$^3$ ozone: CVD mortality
RR 1.0105 (1.0027--1.0185), morbidity RR 1.0578 (1.0289--1.0874). The exposure-window
split is the result: \textbf{warm-season ozone RR 1.0186 (1.0119--1.0253) versus an
annual-average RR of 0.9971 (0.9781--1.0164)} --- annual averaging destroys the signal.
Female RR 1.1135 (1.0057--1.2328) versus male 1.0707 (0.9855--1.1633); Asia estimates are
imprecise. Co-pollutant record; no PM estimate is reported. The averaging-window lesson
transfers directly to PM network design.}
{10.46234/ccdcw2026.161}

\clearpage

\band{Reproductive \& developmental\hfill 1 record}{Coral}

\paperentry{Li et al.\ 2026}
{Prenatal NO$_2$ and stillbirth in 3.98 million Iranian mother--baby pairs}
{Environ Pollut}
{3{,}982{,}304 pairs, 2013--2018, pregnancy-average NO$_2$ from satellite-based gridded
estimates averaged within 10\,km of the delivery hospital, analysed under both the Iranian
($\geq$22 weeks) and WHO ($\geq$28 weeks) stillbirth definitions. Per 10\,$\mu$g/m$^3$:
\textbf{OR 1.206 (1.192--1.220)} and \textbf{1.106 (1.088--1.125)} respectively,
essentially unchanged on adjustment for PM$_{2.5}$ and ozone. The exposure--response is
monotonic with \emph{steeper slope at low concentrations}. A 10\,$\mu$g/m$^3$ target would
avert 7{,}686 (7{,}242--8{,}127) stillbirths, an attributable fraction of 25.8\%.
Exposure assigned at the delivery hospital rather than residence is a substantial
misclassification channel in a country with heavy referral flows, and the effect size is
large enough to warrant scepticism. Co-pollutant, not PM.}
{10.1016/j.envpol.2026.128908}

\band{Occupational \& indoor exposure\hfill 2 records}{Clay}

\paperentry{Makanjuola 2026}
{Personal PM exposure and spirometry among 563 Southwest Nigerian street sweepers}
{Research Square (preprint)}
{Cross-sectional, non-smoking sweepers aged 18--60 from Ado-Ekiti, Ibadan and Ikeja;
personal PM$_{2.5}$ and PM$_{10}$ by handheld PMS5003 at breathing height; spirometry to
ATS standard. Mean PM$_{2.5}$ \textbf{103.55$\pm$72.03\,$\mu$g/m$^3$} and PM$_{10}$
105.09$\pm$73.95, roughly 20$\times$ and 7$\times$ the WHO 2021 annual guideline.
PM$_{2.5}$ correlates inversely with FEV$_1$/FVC ($r=-0.14$, $p=0.001$) but
\textbf{the association is fully attenuated in the adjusted model ($\beta=-0.0001$,
$p=0.286$)}. A bare PMS5003 without RH correction at these loadings is a real
measurement-error problem the paper does not address, and PM$_{2.5}\approx$PM$_{10}$ is
itself a red flag for the instrument. Reported for the exposure characterisation, not the
health inference.}
{10.21203/rs.3.rs-9928829/v1}

\paperentry{Piedra et al.\ 2026}
{Informal employment versus cantonal satellite pollution as respiratory-symptom predictors}
{Front Public Health}
{1{,}807 Ecuadorian workers across four sectors, ECRHS symptom instrument cross-referenced
with Sentinel-5P retrievals; optimal spatial/temporal concentration percentile selected
per outcome by AIC; Firth penalised logistic regression. \textbf{Informal employment
dominates}: asthma aOR 12.56 (5.09--32.87), wheeze 6.78 (1.83--27.89), chronic cough 4.14
(2.26--7.74). After adjusting for cantonal altitude (itself protective, aOR 0.78,
0.66--0.93), NO$_2$ and SO$_2$ remain independent predictors of chronic phlegm and cough
while ozone loses significance. Sentinel-5P carries no PM product, so the ambient arm is
entirely gas-phase; cantonal assignment is coarse and the design is cross-sectional. The
interesting move is using satellite retrievals as an occupational-surveillance covariate.}
{10.3389/fpubh.2026.1874197}

\band{Mechanistic toxicology\hfill 4 records}{Violet}

\paperentry{Silva et al.\ 2026}
{PM$_{2.5}$ organic extracts in 2D monolayers versus 3D spheroids of bronchial epithelium}
{Environ Toxicol Pharmacol}
{16HBE cells in paired 2D and 3D culture exposed to solvent-extractable organic fractions
of well-characterised ambient PM$_{2.5}$ from three sampling periods.
\textbf{Concordance between formats is endpoint-dependent, not uniform}: WST-1 metabolic
activity, ROS fluorescence and extracellular metabolite exchange respond more strongly in
2D, while pro-inflammatory cytokine secretion is more evident in spheroids. Cell death is
necrotic rather than apoptotic in both. The practical implication for in vitro PM
toxicology is that model choice silently selects which hazard you can detect --- 2D favours
direct cytotoxic and oxidative readouts, 3D favours responses that need spatial
organisation and exposure gradients.}
{10.1016/j.etap.2026.105125}

\paperentry{Mishra et al.\ 2026}
{OMA1--DELE1 as a shared node across ultrafine PM, black carbon and nanoplastic immunotoxicity}
{bioRxiv (preprint)}
{Human PBMCs exposed to ultrafine PM, black carbon and polystyrene nanoplastics under
matched conditions, profiled across NRF2, HIF-1$\alpha$, PGC-1$\alpha$, TFAM, OMA1, DELE1,
mitochondrial ND1, complexes I--V, NF-$\kappa$B, TNF-$\alpha$ and NLRP3, with PCA,
hierarchical clustering and correlation networks. All three cause oxidative stress and
mitochondrial dysfunction but with distinct kinetics: \textbf{UFPM gives an acute
antioxidant/adaptive response, black carbon a chronic OMA1--DELE1 integrated stress
response with inflammation, nanoplastics a low-grade persistent adaptation.} The
comparative design under matched dosing is what makes this worth reading; it is a
preprint, expression-level only, with no functional respirometry.}
{10.64898/2026.07.29.741662}

\paperentry{Wang et al.\ 2026}
{Multigenerational reproductive toxicity of three vehicular PM types in \textit{C.\ elegans}}
{Ecotoxicol Environ Saf}
{SRM diesel PM, WHTC diesel PM and WLTP gasoline PM, differing in morphology and elemental
composition, all produced dose-dependent brood-size reduction, germ-cell apoptosis, lower
sperm count and morphological abnormality in parents. Mechanistically: suppressed
SPE-8/SPE-9/major sperm protein/vitellogenin, activated \textit{ced-3/-4/-9}, disrupted
\textit{set-2}, \textit{met-2}, \textit{spr-5}. \textbf{Effects persist to F1--F2 and
recover by F3--F4 --- intergenerational carryover, not stable transgenerational
inheritance}, which is the distinction most papers in this space elide. Nematode model;
extrapolation to mammalian germline is speculative.}
{10.1016/j.ecoenv.2026.120609}

\paperentry{Sharma et al.\ 2026}
{Dim light at night modulates the neuroimmune response to burn-pit emission inhalation}
{J Toxicol Environ Health A}
{32 male Swiss Webster mice housed under dark nights or dim light at night, exposed at
week 4 to filtered air or burn-pit emissions (6.9\,mg/m$^3$ mass,
4.03$\times$10$^6$ particles/cc) for 4\,h/day over 6 days. \textbf{No behavioural
differences} in open field, Y-maze or elevated plus maze. Hippocampal \textit{Il6},
\textit{Tnf}$\alpha$, \textit{Il1}$\beta$ and cortical \textit{Il6} were
\emph{decreased} in LD-exposed mice, with reduced hippocampal \textit{Bmal1}; lung and
liver were unchanged. Small $n$ per cell, a single sex, and a cytokine direction opposite
to the usual inflammatory prior --- reported as a circadian-interaction signal, not as an
inflammation result.}
{10.1080/15287394.2026.2708884}

\clearpage

\band{Burden, policy \& mitigation\hfill 5 records}{Sage}

\paperentry{Yoon et al.\ 2026}
{Heat versus bushfire-smoke PM$_{2.5}$ across four arms of the Australian health system}
{Med J Aust}
{National ED presentations, inpatient admissions, MBS primary-care claims and PBS
prescriptions, 2014--2020, for cardiovascular, respiratory and mental-health conditions,
supporting Australia's first National Climate Risk Assessment. \textbf{Heat loads the
front door} (ED RR 1.033 per $^{\circ}$C for respiratory, 1.023 for mental health) while
\textbf{PM$_{2.5}$ loads admissions and prescribing} (RR 1.014 per 1\,$\mu$g/m$^3$ at
1-day lag for both cardiovascular and respiratory admissions; primary-care mental-health
visits RR 1.049, 1.040--1.058). Mental-health services are sensitive to both exposures
across every domain --- the least-anticipated result and the one with the weakest
mechanistic account. Ecological in exposure assignment.}
{10.5694/mja2.70261}

\paperentry{Slavik et al.\ 2026}
{Objective versus subjective wildfire-smoke exposure and parental protective behaviour}
{Environ Challenges}
{2{,}086 parents in wildfire-prone western US and Canada, surveyed and linked to 3.5 years
of smoke-monitoring data (2020--2023). Perceived exposure (sight or smell) tracks measured
exposure but predicts behaviour far better: \textbf{$b=1.11$ (0.92--1.30) additional
protective actions for perceived exposure versus $b=0.23$ (0.06--0.40) per
10\,$\mu$g/m$^3$ of measured PM$_{2.5}$}. The measured-exposure relationship is
non-monotonic --- rise, plateau, rise. Self-reported behaviour is the obvious weakness and
recall is correlated with perception by construction, which inflates the contrast. Even
discounted, the direction is a design constraint on any public-facing sensor network.}
{10.1016/j.envc.2025.101362}

\paperentry{Wan 2026}
{Place-based ecological restoration and self-rated health: quasi-experimental evidence}
{Econ Hum Biol}
{Six waves of the China Family Panel Studies against the staggered prefecture-level
rollout of China's national ecological restoration pilot, two-way fixed-effects DiD.
\textbf{+0.1697 points of self-rated health on a five-point scale}, no differential
pre-trends in event study, robust to placebo tests, Goodman--Bacon decomposition,
propensity-score matching, double machine learning, IV, and adjustment for overlapping
environmental policies. Mediation runs through lower PM$_{2.5}$, greater greenness and
better visibility alongside economic and psychosocial channels. The unusually thorough
robustness battery is the reason to take it seriously; the mediation analysis is
correlational and the PM$_{2.5}$ share is not separately identified.}
{10.1016/j.ehb.2026.101636}

\paperentry{Zaid et al.\ 2026}
{Fibrous aerosol filtration: theory, CFD assumptions, metrology and fabrication}
{J Mater Sci}
{Review integrating single-fibre capture mechanisms and most-penetrating-particle-size
theory with electrostatic charge generation, transport and stability; 3D reconstruction of
filter geometry; CFD, particle tracking, deposition and clogging simulation; and
electrospinning, meltblowing, bubble electrospinning and supersonic solution blowing. The
section worth reading is the critique of the homogeneity and steady-state assumptions that
most published simulations inherit, alongside an explicit treatment of instrumentation
limits and measurement uncertainty in filtration testing. Review, no new data; the
material--process--performance framing is the organising contribution.}
{10.1007/s10853-026-13269-8}

\paperentry{Sidharaddi et al.\ 2026}
{Electrodeposited copper nodule films as electrostatic particulate collectors}
{Research Square (preprint)}
{Copper nodules electrodeposited onto steel; XRD confirms polycrystalline FCC, SEM a
nodular microstructure, AFM a non-monotonic dependence of grain size (86.5$\to$224.6\,nm)
and roughness (23.4$\to$76.2\,nm) on bath concentration before both fall again. The
0.075\,M film (721\,nm thick) gave \textbf{88.78\% removal of dhoop smoke and an
air-purifying factor of 54.12 over 405\,s}, exceeding a stainless-steel discharge rod.
Chamber test with a single combustion source, no reference-grade particle instrument, no
size resolution and no ozone measurement --- which for a corona-type collector is the
omission that matters. Included as an engineering-control data point, tier C.}
{10.21203/rs.3.rs-10126181/v1}

\band{Other clinical endpoints\hfill 2 records}{Slate}

\paperentry{Jung et al.\ 2026}
{PM$_{2.5}$ and rheumatoid arthritis disease activity and flare}
{Ann Rheum Dis}
{Prospective cohort at a South Korean tertiary centre, 2021--2024: 1{,}070 patients,
12{,}583 outpatient visits, monthly mean concentrations of six pollutants, generalised
estimating equations adjusted for demographics, serology, medication, socioeconomic and
meteorological covariates, plus a within-patient case-crossover sensitivity analysis on
daily concentrations. \textbf{PM$_{2.5}$ alone was associated with flare (adjusted OR
1.113, 95\% CI 1.017--1.218)} and with higher DAS28-CRP, CDAI, and tender and swollen
joint counts; the association was stronger in women and non-smokers, and the
case-crossover arm reproduced it for cumulative exposure beyond two weeks. Single centre,
monitor-assigned exposure; the two-design agreement is what makes it convincing.}
{10.1016/j.ard.2026.06.018}

\paperentry{Lees 2026}
{PM$_{2.5}$ component mixtures and inpatient hospitalisation in chronic kidney disease}
{Research Square (preprint)}
{2{,}511 \textit{All of Us} participants with ICD-coded CKD, 19{,}656 hospitalisations over
22{,}448 person-years (2010--2019); NASA SEDAC estimates of organic and elemental carbon,
ammonium, nitrate and sulfate linked at 3-digit ZIP; weighted quantile sum regression
inside a negative binomial model. Median PM$_{2.5}$ was only 4.4\,$\mu$g/m$^3$.
\textbf{IRR 1.31 (0.99--1.73, $p=0.06$) per quartile of the WQS index --- formally null,
with 93\% of the mixture weight on nitrate}, 4\% each on ammonium and elemental carbon,
and zero on organic carbon and sulfate. WQS weights are unstable at this sample size and
3-digit ZIP assignment is very coarse. The nitrate concentration of weight is the
hypothesis worth carrying, not the point estimate.}
{10.21203/rs.3.rs-10438586/v1}

\clearpage

% ==================================================================== TABLE
\section*{\sffamily\bfseries\color{Deep}Corpus register}
\vspace{-1.5mm}{\color{Amber}\rule{\linewidth}{1.1pt}}\vspace{2.5mm}

{\fontsize{6.7}{8.4}\selectfont
\setlength{\tabcolsep}{2.4pt}
\begin{longtable}{@{}>{\RaggedRight}p{25mm}>{\RaggedRight}p{41mm}>{\RaggedRight}p{27mm}%
>{\RaggedRight}p{23mm}>{\RaggedRight}p{18mm}>{\RaggedRight}p{15mm}@{}}
\toprule
\textsf{\textbf{First author}} & \textsf{\textbf{Focus}} & \textsf{\textbf{Journal}} &
\textsf{\textbf{Design}} & \textsf{\textbf{Metric}} & \textsf{\textbf{Region}}\\
\midrule
\endfirsthead
\toprule
\textsf{\textbf{First author}} & \textsf{\textbf{Focus}} & \textsf{\textbf{Journal}} &
\textsf{\textbf{Design}} & \textsf{\textbf{Metric}} & \textsf{\textbf{Region}}\\
\midrule
\endhead
\multicolumn{6}{@{}l@{}}{\cellcolor{Mist}\textsf{\textbf{\textcolor{Deep}{Neuro / mental health}}}}\\[0.6mm]
\href{https://doi.org/10.1001/jamanetworkopen.2026.25015}{Dales} & Wildfire-sourced PM$_{2.5}$ and ED visits for migraine and other primary headache & \textit{JAMA Netw Open} & Case-crossover & PM$_{2.5}$ (wildfire) & Canada\\
\href{https://doi.org/10.1093/ageing/afag215}{Glover} & Umbrella meta-review of urban environmental exposures and cognitive health & \textit{Age Ageing} & Umbrella review / meta-analysis & PM$_{2.5}$ & Global\\
\href{https://doi.org/10.1016/j.nsa.2026.107019}{Goossens} & In utero air pollution exposure and pre-/postnatal brain development & \textit{Neurosci Appl} & Narrative review & PM$_{2.5}$ / PM$_{10}$ / PM$_{\mathrm{coarse}}$ & Global\\
\href{https://doi.org/10.1016/j.isci.2026.116617}{Ma} & PM$_{2.5}$-induced neurotoxicity and depression: exposure risks and mechanisms & \textit{iScience} & Narrative review & PM$_{2.5}$ + components & Global\\
\href{https://doi.org/10.3760/cma.j.cn112150-20250804-00752}{Pan} & Maternal PM exposure and offspring neurodevelopment: effects and mechanisms & \textit{Zhonghua Yu Fang Yi Xue Za Zhi} & Narrative review & PM (bulk) & China\\
\addlinespace[0.9mm]
\multicolumn{6}{@{}l@{}}{\cellcolor{Mist}\textsf{\textbf{\textcolor{Teal}{Cardiovascular \& metabolic}}}}\\[0.6mm]
\href{https://doi.org/10.3760/cma.j.cn112338-20251120-00843}{Huang} & PM$_{2.5}$ and dyslipidaemia in rural Chinese children; dietary effect modification & \textit{Zhonghua Liu Xing Bing Xue Za Zhi} & Cross-sectional (multistage) & PM$_{2.5}$ & China\\
\href{https://doi.org/10.32394/pe/218966}{Kazmierska} & Air pollution and cardiovascular disease: review of current data & \textit{Przegl Epidemiol} & Narrative review & PM$_{2.5}$ / PM$_{10}$ & Global\\
\href{https://doi.org/10.1007/s00125-026-06814-2}{Xi} & Air pollution, built environment and temperature vs. prediabetes phenotypes & \textit{Diabetologia} & Prospective cohort & PM$_{2.5}$ & Germany\\
\href{https://doi.org/10.1161/JAHA.125.046917}{Yang} & Short-term PM$_{2.5}$ and hospitalisation/ED use in atrial fibrillation & \textit{J Am Heart Assoc} & Case-crossover & PM$_{2.5}$ (wildfire-influenced) & USA\\
\addlinespace[0.9mm]
\multicolumn{6}{@{}l@{}}{\cellcolor{Mist}\textsf{\textbf{\textcolor{Sky}{Reproductive \& developmental}}}}\\[0.6mm]
\href{https://doi.org/10.1016/j.jes.2025.10.043}{Chen} & Prenatal air pollution, fetal growth trajectory phenotypes and child neurodevelopment & \textit{J Environ Sci (China)} & Prospective cohort & PM$_{2.5}$ / PM$_{10}$ & China\\
\href{https://doi.org/10.1016/j.ecoenv.2026.120566}{Huang} & Joint prenatal bisphenol + air-pollutant mixture exposure and preschool IQ & \textit{Ecotoxicol Environ Saf} & Prospective cohort (mixtures) & PM$_{2.5}$ / PM$_{10}$ & China\\
\href{https://doi.org/10.1007/s40572-026-00552-8}{Peterson} & Ambient air pollution, PAHs and pubertal development: critical review & \textit{Curr Environ Health Rep} & Critical review & PM$_{2.5}$ components & Global\\
\href{https://doi.org/10.1016/j.envres.2026.125313}{Xie} & Long-term PM$_{2.5}$ and timing of pubertal development in boys and girls & \textit{Environ Res} & Prospective cohort & PM$_{2.5}$ & China\\
\addlinespace[0.9mm]
\multicolumn{6}{@{}l@{}}{\cellcolor{Mist}\textsf{\textbf{\textcolor{Sage}{Respiratory \& allergic}}}}\\[0.6mm]
\href{https://doi.org/10.1007/s00484-026-03289-z}{Bayar Muluk} & Climate change and otolaryngology: airway disease impacts and emission reduction & \textit{Int J Biometeorol} & Narrative review & PM$_{2.5}$ & Global\\
\href{https://doi.org/10.4103/mgr.MEDGASRES-D-25-00164}{Liao} & Short-term air pollution, respiratory symptoms and lung function in asthma & \textit{Med Gas Res} & Panel study & PM$_{2.5}$ / PM$_{10}$ & China\\
\href{https://doi.org/10.3390/epidemiologia7040097}{Sanchez de Prada} & Ambient PM and asthma-coded admissions in children 0-3 y, Spain & \textit{Epidemiologia} & Ecological panel & PM$_{2.5}$ / PM$_{10}$ & Spain\\
\href{https://doi.org/10.1016/j.jhazmat.2026.143076}{Yang} & Allergenic vs. toxic risk of airborne seed hairs compared with ambient PM$_{2.5}$ & \textit{J Hazard Mater} & Field measurement & PM$_{2.5}$ + bioaerosol & China\\
\addlinespace[0.9mm]
\multicolumn{6}{@{}l@{}}{\cellcolor{Mist}\textsf{\textbf{\textcolor{Violet}{Mechanistic toxicology}}}}\\[0.6mm]
\href{https://doi.org/10.1016/j.jes.2025.12.074}{Ding} & PM$_{2.5}$ promotes vascular calcification via miR-27b-5p/HES1 signalling & \textit{J Environ Sci (China)} & In vitro + organoid + murine & PM$_{2.5}$ & China\\
\href{https://doi.org/10.1007/s00360-026-01698-5}{Maraschi} & Osmoregulatory disruption and metal bioaccumulation from settleable atmospheric PM & \textit{J Comp Physiol B} & Controlled exposure (ecotox) & Settleable APM & Brazil\\
\href{https://doi.org/10.1016/j.jes.2025.11.018}{Wu} & Macrophage NLRP3 inflammasome in real-ambient PM-induced hepatic glucose dysregulation & \textit{J Environ Sci (China)} & Real-ambient murine + co-culture & Real-ambient PM & China\\
\addlinespace[0.9mm]
\multicolumn{6}{@{}l@{}}{\cellcolor{Mist}\textsf{\textbf{\textcolor{Clay}{Occupational \& indoor}}}}\\[0.6mm]
\href{https://doi.org/10.3389/fmedt.2026.1796828}{Chen} & Nested suction instrument to cut intraoperative surgical-smoke PM$_{2.5}$ & \textit{Front Med Technol} & Bench / device experiment & PM$_{2.5}$ (surgical smoke) & China\\
\href{https://doi.org/10.1093/annweh/wxag055}{Cole} & Determinants of respirable crystalline silica exposure in Australian tunnelling & \textit{Ann Work Expo Health} & Observational exposure study & Respirable dust / RCS & Australia\\
\href{https://doi.org/10.1016/j.jhazmat.2026.143081}{Ding} & Stage-resolved hair Zn-isotope dynamics as a firefighting combustion-exposure tracer & \textit{J Hazard Mater} & Repeated-measures biomarker & Combustion-derived PM & China\\
\href{https://doi.org/10.3390/jox16040119}{Ramsden} & Xenobiotic hazards in aircraft cabin air & \textit{J Xenobiot} & Review + risk economics & Ultrafine particles & Global\\
\addlinespace[0.9mm]
\multicolumn{6}{@{}l@{}}{\cellcolor{Mist}\textsf{\textbf{\textcolor{Amber}{Sensing, forecasting \& instrumentation}}}}\\[0.6mm]
\href{https://doi.org/10.1016/j.jenvman.2026.130515}{Shi} & Adaptive spatiotemporal graph transformer for multi-site multi-step PM$_{2.5}$ forecasting & \textit{J Environ Manage} & Deep-learning model & PM$_{2.5}$ & China / India\\
\href{https://doi.org/10.3390/bios16070380}{Tsuchiya} & Seasonal and indoor-outdoor characteristics of fluorescent aerosol particles in offices & \textit{Biosensors} & Real-time sensor campaign & Bioaerosol / coarse & Japan\\
\addlinespace[0.9mm]
\multicolumn{6}{@{}l@{}}{\cellcolor{Mist}\textsf{\textbf{\textcolor{Amber}{Exposure assessment \& modelling}}}}\\[0.6mm]
\href{https://doi.org/10.1029/2025GH001621}{Solorzano-Araque} & Multiscale exposure modelling: LOTOS-EUROS + mobility surveys + particle toxicity & \textit{GeoHealth} & CTM + mobility integration & PM$_{2.5}$ / PM$_{10}$ & Colombia\\
\addlinespace[0.9mm]
\multicolumn{6}{@{}l@{}}{\cellcolor{Mist}\textsf{\textbf{\textcolor{Coral}{Burden, policy \& mitigation}}}}\\[0.6mm]
\href{https://doi.org/10.1002/hsr2.72886}{Karamzad} & Ischaemic heart disease burden in MENA 1990-2021; PM attributable fraction 28.2\% & \textit{Health Sci Rep} & GBD secondary analysis & Ambient PM & MENA\\
\href{https://doi.org/10.1371/journal.pone.0347989}{Rees} & HEPA filtration to reduce respiratory infection in care homes & \textit{PLoS One} & RCT process evaluation & Filtered PM & UK\\
\href{https://doi.org/10.1016/j.jes.2026.01.009}{Shan} & Policy-driven emission reductions and post-pandemic rebound, high-resolution inventories & \textit{J Environ Sci (China)} & Emission inventory & PM emissions & China\\
\addlinespace[0.9mm]
\multicolumn{6}{@{}l@{}}{\cellcolor{Mist}\textsf{\textbf{\textcolor{Slate}{Other clinical endpoints}}}}\\[0.6mm]
\href{https://doi.org/10.3390/dj14070440}{Lupita} & PM$_{2.5}$, preventive-care capacity and caries/edentulism burden across the EU & \textit{Dent J} & Ecological panel & PM$_{2.5}$ & EU-27\\
\href{https://doi.org/10.1016/j.jdin.2026.06.008}{Park} & Long-term PM exposure and risk of androgenetic alopecia & \textit{JAAD Int} & Nationwide cohort & PM$_{2.5}$ / PM$_{10}$ & South Korea\\
\href{https://doi.org/10.3171/2026.3.PEDS2615}{Sun} & Environmental exposure and area deprivation vs. paediatric hydrocephalus outcomes & \textit{J Neurosurg Pediatr} & Registry cohort + geospatial & PM$_{2.5}$ (satellite-derived) & USA\\
\addlinespace[0.9mm]
\bottomrule
\end{longtable}}


\vspace{4mm}
\begin{keybox}
{\sffamily\bfseries\footnotesize\color{Deep}Provenance and method}\\[1.2mm]
{\fontsize{7.2}{8.8}\selectfont
Entry window \textbf{6 August 2026}, contiguous with the 5 August issue; no gap. Most
shipped records carry a 6 August PubMed entry date, Europe PMC creation date or Crossref
\texttt{created} date; five PubMed records carry a 4--5 August entry date and are
\textbf{explicitly flagged backfill} --- PubMed re-indexed them into today's window and
\texttt{seen.json} confirms they were never shipped. They must be excluded from any
weekly trend line.
\textbf{Sources.} PubMed E-utilities and the PubMed connector on two separate axes
(health, instrumentation) windowed on \texttt{[EDAT]}: the connector returned 20 health
and 4 sensing hits, the local harvester 14 and 4, and \textbf{the two disagree in both
directions} --- ten records were connector-only and three harvester-only, so both legs are
merged before screening. Union 25 records, 18 carried, 7 rejected (health-literacy
instruments, e-cigarette nanoparticles, cigarette-smoke miRNA, marine particulate organic
carbon, dredged-sludge stabilisation, a COPD protease review with no PM exposure, and one
duplicate already shipped on 5 August). Europe PMC on both \texttt{FIRST\_PDATE} and
\texttt{CREATION\_DATE} --- 7 preprints, 5 carried, 1 rejected as off-topic
(GPU lattice-Boltzmann radiative transfer), 1 dropped as already shipped. Crossref scoped
by journal ISSN across eight aerosol and atmospheric titles PubMed does not index ---
11 hits, 4 carried, 5 rejected out of scope, 2 held for want of an abstract. Consensus on
three sensing axes (calibration and drift, satellite--surface fusion, humidity correction
and edge inference) --- \textbf{all nine hits resolved to DOIs created between January and
June 2026}, i.e.\ entirely outside this window, and were rejected rather than backfilled;
Ginsburg et al.'s global PurpleAir calibration is flagged in the run log as the priority
pickup if a backfill slot opens. OpenAlex returned \texttt{HTTP 429} on top of the
existing paid block. arXiv returned 60 records of which 2 postdate 25 July, neither in
scope after the recency screen. Scholar Gateway was queried for the held Prajapati
hygroscopicity record and returned adjacent literature only --- its corpus is current to
May 2026.
\textbf{Screening.} Deduplication keys PMID $\rightarrow$ lowercased DOI $\rightarrow$
title signature, checked against \texttt{seen.json} \emph{before} summarising. Consensus
titles were resolved to DOIs through Crossref before screening. \textbf{Nine records are
held, not rejected}: DOIs verify in Crossref but no abstract is retrievable from Crossref,
Europe PMC, Semantic Scholar or Consensus. Seven are carried over from 4--5 August and
were re-tried on this run without success (a third attempt for four of them); two are new
Elsevier deposits. Held records stay out of \texttt{seen.json} so a later issue can carry
them. Nothing was summarised from a title alone.
\textbf{Labels.} Subtopic, design, particle-metric and geography assignments are
single-label judgements from title, abstract, keywords and MeSH terms, not a controlled
vocabulary. The six effect estimates are transcribed verbatim; two are co-pollutant
(NO$_2$, O$_3$) rather than PM and are marked as such in the figure caption and the
register.
All DOI links resolve to the publisher of record and passed the automated DOI gate
(0 fail, 7 advisory title-overlap warnings, all from editorial focus lines) before this
issue compiled. Abstracts remain the copyright of their publishers.\par}
\end{keybox}

\end{document}
